\subsection{Wizualizacja t-SNE słów}\label{sec:we-tsne}

Rys.~\ref{fig:we-tsne-ipi} i~\ref{fig:we-tsne-pol} pokazują t-SNE (odległość kosinusowa) wszystkich słów z wektorem.
Wersje interaktywne -- po najechaniu na punkt widać słowo, jego etykietę, liczbę anotacji każdej etykiety i rozkład
etykiet zdań; lista rozwijana zmienia kolorowanie -- są w repozytorium
(\texttt{results/word\_embeddings/tsne\_words\_*.html}). Ponieważ t-SNE zniekształca odległości, grupowanie
sprawdzono też w oryginalnej przestrzeni (tab.~\ref{tab:we-separation}).

\begin{figure}[H]
\centering
\includegraphics[width=0.95\textwidth]{tsne_words_w2v_all.png}
\caption{t-SNE słów -- word2vec IPI PAN. Kolor: etykieta słowa (lewy) i dominująca etykieta zdań, w których słowo
wystąpiło (prawy); wielkość punktu rośnie z liczbą wystąpień, słowa bez etykiety są szarym tłem.}
\label{fig:we-tsne-ipi}
\end{figure}
\begin{figure}[H]
\centering
\includegraphics[width=0.95\textwidth]{tsne_words_w2v_polemo_all.png}
\caption{t-SNE słów -- word2vec PolEmo. Kolor: etykieta słowa (lewy) i dominująca etykieta zdań, w których słowo
wystąpiło (prawy); wielkość punktu rośnie z liczbą wystąpień, słowa bez etykiety są szarym tłem.}
\label{fig:we-tsne-pol}
\end{figure}

\begin{table}[H]
\centering\small
\begin{tabular}{llrrrrrr}
\toprule
Model & Etykiety & $n$ & czystość k-NN & losowo & sylwetka & sonda & losowo \\
\midrule
IPI PAN & WYR\_POZ / WYR\_NEG & 215 & 0{,}648 & 0{,}536 & 0{,}033 & 0{,}72 & 0{,}50 \\
IPI PAN & POZ / NEG (zdania) & 622 & 0{,}545 & 0{,}502 & 0{,}004 & 0{,}62 & 0{,}50 \\
PolEmo & WYR\_POZ / WYR\_NEG & 181 & 0{,}671 & 0{,}571 & 0{,}038 & 0{,}58 & 0{,}50 \\
PolEmo & POZ / NEG (zdania) & 569 & 0{,}559 & 0{,}504 & 0{,}006 & 0{,}62 & 0{,}50 \\
\bottomrule
\end{tabular}
\caption{Czy etykiety grupują się w przestrzeni osadzeń (100 wymiarów, kosinus). Czystość k-NN -- średni odsetek
10 najbliższych sąsiadów z tą samą etykietą; „losowo'' -- ta sama miara po przetasowaniu etykiet (1000 razy, $p<0{,}01$
we wszystkich wierszach). Sonda -- zrównoważona trafność regresji logistycznej (5-krotna walidacja krzyżowa).}
\label{tab:we-separation}
\end{table}
